\chapter{Values and Variables}\label{ch-values}

A program works on \emph{values}: numbers, pieces of text, true and
false, lists of things. This chapter is about the values Zeno has, and
about the names you give them.

\section{The kinds of value}

Every value has a \emph{type}, which says what kind of value it is and
what can be done with it. These are the types every program meets
first:

\begin{center}
\begin{tabular}{lll}
\toprule
Type & What it holds & Written like \\
\midrule
\verb|int| & a whole number & \verb|42|, \verb|-7|, \verb|0xFF|, \verb|1_000_000| \\
\verb|float| & a number with a fraction & \verb|3.14|, \verb|-0.5|, \verb|2.5e-7| \\
\verb|bool| & true or false & \verb|true|, \verb|false| \\
\verb|str| & text & \verb|"Ada"| \\
\verb|list[int]| & several values, in order & \verb|[1, 2, 3]| \\
\verb|dict[str, int]| & values filed under keys & \verb|["Ada": 36]| \\
\verb|none| & no value & \verb|none| \\
\bottomrule
\end{tabular}
\end{center}

\begin{zeno}
print(42, 3.14, true, "Ada", [1, 2, 3], ["Ada": 36], none)
\end{zeno}

Later chapters add more: numbers with a unit (\verb|500 ms|), numbers
with a size (\verb|u8|), objects of your own classes, enum values, and
functions.

\section{let: a name for a value}

\verb|let| gives a value a name:

\begin{zeno}
let name = "Ada"
let born = 1815
print(name, "was born in", born)
\end{zeno}

A name given with \verb|let| keeps its value for as long as it exists.
Trying to give it another is a mistake, and the message names the word
that would allow it:

\begin{zeno}[fails=yes]
let count = 0
count = 1
\end{zeno}

\section{var: a name whose value changes}

When the value is meant to change, declare the name with \verb|var|.
After that, plain \verb|=| gives it a new value:

\begin{zeno}
var count = 0
count = count + 1
count = count + 1
print(count)
\end{zeno}

There is a shorter way to write \verb|count = count + 1|:
\verb|count += 1|. Every operator has a form like it, and
\hyperref[ch-numbers]{Numbers and Operators} lists them.

A name is declared once, and only declared names can be given a value.
So a misspelt name is a mistake, and never a new variable by accident:

\begin{zeno}[fails=yes]
var count = 0
cuont = 1
\end{zeno}

\begin{why}
In Python, \verb|x = 1| both makes a new variable and changes an old
one, so a typing slip quietly makes a second variable and the first
one never changes. Separating the two turns the slip into an error.
\verb|let| is the short, ordinary word because most values never
change once they are made, and a program whose values do not change
is easier to follow. (Decision 1.6.)
\end{why}

\section{const: a value fixed for the whole program}

\begin{zeno}
const LIMIT = 100
const GREETING = "hello"
print(LIMIT, GREETING)
\end{zeno}

A \verb|const| says that the value belongs to the program, not to a
moment while it runs: a size, a limit, an address, a name. By custom
its name is written in capitals. Like a \verb|let|, it can never be
given another value.

\section{Types are optional}

You can write the type of a name after a \verb|:|, or leave it out:

\begin{zeno}
let port: int = 8080
let host = "localhost"          # a str: Zeno sees that
var ratio: float = 0.5
print(host, port, ratio)
\end{zeno}

Where the type is left out, Zeno works it out from the value. Where it
is written, the value must fit it, and from then on every value given
to the name must fit it too:

\begin{zeno}[fails=yes]
var level: int = 3
level = "high"
\end{zeno}

A whole number is not quietly made a float, even where a float is
wanted. Write the decimal point:

\begin{zeno}[fails=yes]
let ratio: float = 1
\end{zeno}

\begin{why}
Turning an \verb|int| into a \verb|float| where one is stored would
leave a value whose written type and real type disagree, and in a
compiled program it hides an instruction. Arithmetic is different:
\verb|1 + 0.5| makes a new value, and nothing is disguised.
(Decisions 2.3 and 2.3b.)
\end{why}

A name must have a value from the moment it is declared. There is no
\verb|var x| with nothing after it:

\begin{zeno}[fails=yes]
var total: int
\end{zeno}

\hyperref[ch-types]{Types, Written and Worked Out} goes into all of
this in depth: what Zeno works out, what it cannot, and the strict
mode in which every type must be known.

\section{Names}

A name is made of letters, digits and \verb|_|, and does not start
with a digit. Letters of any alphabet are allowed. Case counts:
\verb|total| and \verb|Total| are two names. By custom, variables and
functions are written \verb|like_this|, classes and enums
\verb|LikeThis|, and constants \verb|LIKE_THIS|.

The words of the language, such as \verb|let|, \verb|if| and
\verb|take|, can never be names, and the message says where the word
goes instead:

\begin{zeno}[fails=yes]
let take = 3
\end{zeno}

The names of types (\verb|int|, \verb|str|, \verb|list|) and of the
built-in functions (\verb|print|, \verb|len|) are not reserved. They
are names like any other, and a program may declare its own. That is
rarely a good idea, since the built-in one is then out of reach in
that block.

\section{Blocks and scope}

A pair of braces \verb|{ }| makes a \emph{block}, and a name declared
in a block exists from its declaration to the end of that block. This
stretch of the program is the name's \emph{scope}:

\begin{zeno}
let x = 1
if true {
    let x = 2          # a new x, for this block only
    let y = 3
    print(x, y)
}
print(x)               # the first x again
\end{zeno}

After the block, \verb|y| is gone:

\begin{zeno}[fails=yes]
if true {
    let y = 3
}
print(y)
\end{zeno}

An inner block may declare a name that an outer block already has,
and inside the inner block the name means the inner one. Two
declarations of one name in the same block are a mistake:

\begin{zeno}[fails=yes]
let speed = 10
let speed = 20
\end{zeno}

\begin{why}
A variable that lives exactly as long as its block has a lifetime you
can see in the layout of the code. In Python a variable made inside an
\verb|if| goes on existing after it. (Decisions 3.1 and 3.2.)
\end{why}

\section{Two names for one list}

A name gives a value a name; it does not make a copy. When the value
is a list or an object, two names can stand for the same one, and a
change made through either is seen through both:

\begin{zeno}
let a = [1, 2]
let b = a
b.append(3)
print(a)
\end{zeno}

This also means that \verb|let| fixes the \emph{name}, not what is
inside the value. \verb|a| above can never be given another list, but
the list it names can still grow. Numbers, text and true or false
cannot be changed inside, so for them the difference never shows.
\hyperref[ch-lists]{Lists} has more on this.

\begin{summary}
Every value has a type. \verb|let| names a value that stays,
\verb|var| one that changes, and \verb|const| one fixed for the whole
program. Plain \verb|=| only changes a name that already exists. Types
can be written after \verb|:| or left for Zeno to work out. A name
lives until the end of the block it was declared in.
\end{summary}

\begin{tryit}
\begin{enumerate}
\item Declare \verb|var price = 10| and then give it \verb|10.5|. What
  does Zeno say, and how would you declare it so that both are
  allowed?
\item Write \verb|let x: str = 5| and \verb|let y: int = "5"|. Read
  both messages.
\item Use \verb|print| as a variable name with \verb|let print = 3|,
  and then try to print something below it. What happened?
\item Make two names for one list, change it through one, and print
  it through the other. Then do the same with two numbers and
  \verb|var|. Why is the result different?
\end{enumerate}
\end{tryit}
